\begin{frame}[fragile]{ES6 Class Feature}

\begin{itemize}
	\item JavaScript ES6 is the version that actively adopts OOP ideas. 
	\item In ES6, JavaScript introduces the `class' concept to support the OOP paradigm. 
	\item However, behind the scene, it is syntactic sugar, meaning that it mimics OOP programming methodology, not that ES6 is an OOP language. 
\end{itemize}


\end{frame}

\begin{frame}[fragile]{}

\begin{itemize}
	\item This example shows how we can use `class' to incorporate the constructor function. 
\end{itemize}

\begin{Code}{}%
\begin{lstlisting}[firstnumber=1, label=glabels, xleftmargin=0pt, basicstyle=\tiny]% 

// ES5 or before
function Machine3(name){
  this.name = name;
  this.hello = function() {
    return('Hello ' + this.name)
  }
}
// ES6
class Machine4 {
  constructor(name) {
    this.name = name
    this.hello = function() {
      return('Hello ' + this.name)
    }
  }
}
\end{lstlisting}%   
\end{Code}%

\end{frame}

\begin{frame}[fragile]{}

When we need to add more prototype functions, we have two approaches.

\begin{itemize}
	\item The first approach is to use a prototype to add a new function (lines 1 - 4).
	\item The second approach is to make the prototype function in the class (lines 14 - 16). 
\end{itemize}

\begin{Code}{}%
\begin{lstlisting}[firstnumber=1, label=glabels, xleftmargin=0pt, basicstyle=\tiny]% 

// Adding a function to prototype - method 1
Machine4.prototype.hello2 = function() {
  return('Hello2 ' + this.name)
}

// Adding a function to prototype - method 2
class Machine5 {
  constructor(name) {
    this.name = name
    this.hello = function() {
      return('Hello ' + this.name)
    }
  }
  hello2() { // added to prototype
    return('Hello2 ' + this.name)
  }
}
\end{lstlisting}%   
\end{Code}%

\end{frame}

\begin{frame}[fragile]{}

\begin{itemize}
	\item We can make an object (line 1). 
	\item In this example, person1 is an object from the class NewMachine5, so it has all the functions and variables in the constructor.
	\item However, even though it does not have hello2() function, it can access and use it (line 3). 
\end{itemize}

\begin{Code}{}%
\begin{lstlisting}[firstnumber=1, label=glabels, xleftmargin=0pt, basicstyle=\scriptsize]% 

var person1 = new NewMachine5('Sam','GH565');
console.log(person1); // Machine5 { name: 'Sam', hello: ...}
console.log(person1.hello2()) // Hello2 Sam
\end{lstlisting}%   
\end{Code}%

\end{frame}

\begin{frame}[fragile]{Getter and Setter}

\begin{itemize}
	\item In CSC 260, we learned the OOP idea, encapsulation, to hide data from direct access. 
	\item From this idea, we can make the function setAge() to access and modify the variable (data) in an object. 
\end{itemize}

\begin{Code}{}%
\begin{lstlisting}[firstnumber=1, label=glabels, xleftmargin=0pt, basicstyle=\scriptsize]% 

var person = {
  name: 'Sam',
  age: 10,
  setAge(age) {
    this.age = parseInt(age); // if age is string type value
  }
}
person.setAge(15)
console.log(person)
\end{lstlisting}%   
\end{Code}%

\end{frame}

\begin{frame}[fragile]{}

\begin{itemize}
	\item With JavaScript, we can use the set/get modifier keyword to make the functions, setAge/nextAge, properties. 
	\item Properties are functions that can be used as if they are variables. 
\end{itemize}

\begin{Code}{}%
\begin{lstlisting}[firstnumber=1, label=glabels, xleftmargin=0pt, basicstyle=\scriptsize]% 

var person = {
  name: 'Sam',
  age: 10,
  set setAge(age) { // set makes setAge function a setAge property
    this.age = parseInt(age); // if age is string type value
  },
  get nextAge() {
    return this.age + 1
  }
}
person.setAge = 20 // setAge becomes a property
console.log(person.nextAge)
console.log(person)
\end{lstlisting}%   
\end{Code}%


\end{frame}